\begin{textsample}{POS Dim 3 – ac – Score 43.00 – ac/ac\_inf\_19.txt}  \label{ex:f3_pos_025}
19. \textbf{CRIANÇAS} BEM \textbf{PEQUENAS} ( 1 ANO E 7 \textbf{MESES} A 3 ANOS E 11 \textbf{MESES} ) 19.1.
Campo de experiências “ Traços, \textbf{sons}, cores e formas ” Conviver com diferentes manifestações artísticas, culturais e científicas, locais e universais, no cotidiano da instituição escolar, possibilita às \textbf{crianças}, por meio de experiências diversificadas, vivenciar diversas formas de expressão e linguagens, como as artes visuais ( \textbf{pintura}, modelagem, colagem, fotografia etc. ), a música, o teatro, a dança e o audiovisual, entre outras.
Com base nessas experiências, elas se expressam por várias linguagens, criando suas próprias produções artísticas ou culturais, exercitando a autoria ( coletiva e individual ) com \textbf{sons}, traços, \textbf{gestos}, danças, \textbf{mímicas}, encenações, canções, desenhos, modelagens, manipulação de diversos materiais e de recursos tecnológicos.
Essas experiências contribuem para que, desde muito \textbf{pequenas}, as \textbf{crianças} desenvolvam senso estético e crítico, o conhecimento de si mesmas, dos outros e da realidade que as cerca.
Portanto, a Educação \textbf{Infantil} precisa promover a participação das \textbf{crianças} em tempos e espaços para a produção, manifestação e apreciação artística, de modo a favorecer o desenvolvimento da sensibilidade, da criatividade e da expressão pessoal das \textbf{crianças}, permitindo que se apropriem e reconfigurem, permanentemente, a cultura e potencializem suas singularidades, ao ampliar repertórios e interpretar suas experiências e vivências artísticas. 19.2.
Direitos de aprendizagem e desenvolvimento - Conviver com outras \textbf{crianças} e \textbf{adultos} em situações prazerosas de artes \textbf{plásticas}, músicas, danças, teatro, cinema e festas da cultura popular nacional e regional, que possibilitem produzir conhecimentos sobre si e sobre o mundo social, físico e cultural. - \textbf{Brincar}, utilizando \textbf{criativamente} o repertório da cultura nas suas diferentes manifestações, que sejam : \textbf{plástica}, visual, sonora e musical, em contextos que possibilitem reflexões sobre sua história e de seu grupo. - Explorar, guiados pela \textbf{curiosidade} e pelo \textbf{desejo} de aprender, instrumentos e objetos que produzam \textbf{sons}, ritmos variados, escultura, desenhos, ilustrações, \textbf{pinturas} e gravuras em diferentes suportes. - Participar de experiências \textbf{lúdicas} e artísticas que desenvolvam posturas diferenciadas quanto à organização de ambiente ( tanto do cotidiano quanto o preparo para determinados eventos ), à definição de temas e à escolha de materiais a serem usados. - Expressar suas ideias, pensamentos e sensações, atribuindo sentido ao mundo, cantando, dançando, esculpindo, desenhando e encenando. - Conhecer -se mediante as descobertas realizadas no contato com as diferentes manifestações artísticas e culturais locais e de outras comunidades. 19.2.1.
Objetivo de aprendizagem e desenvolvimento : ( EI02TS01 ) - Criar \textbf{sons} com materiais, objetos e instrumentos musicais, para acompanhar diversos ritmos de música.
Aprendizagens específicas - Expressar -se e interagir com seus pares e \textbf{adultos}, explorando os elementos da música ( melodia, ritmo e harmonia ). - Desenvolver a percepção auditiva através das diferentes possibilidades sonoras. - Desenvolver crescente habilidade de concentração e atenção, utilizando diferentes instrumentos e objetos para reproduzir canções. - Reproduzir ou \textbf{inventar} canções a partir de outras conhecidas. - Desenvolver, \textbf{progressivamente}, a percepção dos \textbf{sons} do seu corpo, dos diversos seres, dos elementos da natureza e de objetos. - Desenvolver a capacidade de interpretação musical envolvendo o corpo no compasso e ritmo da canção. - Acompanhar a música utilizando diferentes objetos ou instrumentos, buscando adequar ao ritmo da música. - Desenvolver a postura investigativa, explorando os \textbf{sons} dos instrumentos musicais e de outros objetos. - Apreciar e valorizar diferentes estilos musicais. - Construir seu gosto musical e fazer escolhas, \textbf{demonstrando} sua preferência.
Sugestões de experiências - \textbf{Brincar} com o próprio corpo \textbf{batendo} \textbf{palmas} e \textbf{pés}. - \textbf{Brincar} com materiais, objetos e instrumentos musicais. - \textbf{Imitar}, \textbf{inventar} e reproduzir criações musicais. - Explorar materiais buscando diferentes \textbf{sons} para acompanhar canções que lhes são familiares. - \textbf{Imitar} \textbf{sons} com diferentes objetos ou instrumentos ao ritmo da música. - \textbf{Brincar} com objetos sonoros e instrumentos musicais de sucatas e/ou industrializados. - Apreciar músicas de produção artística da cultura local e de outras culturas. - \textbf{Brincar} de fazer vozes diferentes, mais finas, mais grossas, mais nasais, dentre outras. - Soprar, sacudir, raspar, \textbf{bater}, experimentando diferentes modos de ação instrumental. - \textbf{Emitir} \textbf{sons} articulados a \textbf{gestos} e ser interpretado pelo outro. - \textbf{Imitar} e \textbf{inventar} \textbf{sons} com o corpo e com objetos diversos. - Construir \textbf{brinquedos} sonoros e outros objetos. 19.2.2.
Objetivo de aprendizagem e desenvolvimento : ( EI02TS02 ) - Utilizar materiais variados com possibilidades de manipulação ( argila, massa de modelar ), explorando cores, \textbf{texturas}, superfícies, planos, formas e volumes ao criar objetos tridimensionais.
Aprendizagens específicas - Expressar -se em diferentes modalidades da linguagem visual – desenho, \textbf{pintura}, colagem, construção, modelagem, utilizando materiais tradicionais e materiais recicláveis e \textbf{reutilizáveis}. - Ampliar as possibilidades de \textbf{manuseio} de materiais diversos para realizar suas produções artísticas. - Aprimorar as habilidades ( liberdade dos \textbf{gestos}, movimento amplo e exploração da dimensão espacial ) necessárias para \textbf{manuseio} de diferentes materiais e instrumentos. - Ampliar conhecimento de mundo a partir da apreciação e participação em diferentes manifestações artísticas e culturais locais e de outras comunidades. - Desenvolver, \textbf{progressivamente}, a sensibilidade artística e a apreciação estética. - Construir atitude de respeito pela produção dos \textbf{colegas} e \textbf{autoconfiança} por sua produção artística. - Desenvolver atitudes de \textbf{cuidado} com o próprio corpo e do \textbf{colega} ao \textbf{manusear} diferentes materiais, instrumentos e objetos durante as produções artísticas. - Desenvolver atitude de colaboração quanto a organização e \textbf{cuidados} com os materiais, suas produções e de seus \textbf{colegas}.
Sugestões de experiências - Criar objetos tridimensionais com palitos de madeira, papéis diversos, argila, massa de modelar, \textbf{areia} de modelar, barro, \textbf{tinta} e outros. - \textbf{Brincar} de \textbf{montar}, encaixar e \textbf{empilhar} objetos tridimensionais de diferentes \textbf{texturas}, formas, \textbf{pesos}, tamanhos, largura e altura. - Pintar e colar formas tridimensionais. - Apreciar as próprias produções e as dos \textbf{colegas}. - \textbf{Manipular} diferentes materiais para desenhar, pintar e modelar. - Observar as transformações dos materiais e as diferentes \textbf{texturas} de massas e \textbf{tintas} produzidas com elementos da natureza. - Experimentar diferentes formas de utilização de materiais, meios e suportes tanto no formato como no tamanho. - Desenhar a partir de uma imagem. - Criar produções artísticas em diferentes suportes ( papel, cartão, cartolina, papelão, madeira, tecido, folhas ou fibras naturais ). - Desenhar a figura humana : a si mesmo, os \textbf{colegas}, o professor, a família. - Produzir colagens utilizando diferentes tipos de cola, materiais e suportes. - Modelar usando diferentes tipos de materiais, com ou sem auxílio de ferramentas. - \textbf{Brincar} com \textbf{tintas} de cores variadas. - Pesquisar e criar diferentes tipos de \textbf{texturas}. - Ter cuidado com o próprio corpo e o dos \textbf{colegas} ao \textbf{manusear} diferentes materiais, instrumentos e objetos. - Cuidar dos materiais e dos trabalhos e objetos produzidos pelo grupo. - Organizar e cuidar do espaço físico. - Apreciar as próprias produções e as produções de outras \textbf{crianças}. - Observar e apreciar produções artísticas e obras de arte. - Criar cenários e figurinos. - Visitar diferentes espaços culturais. 19.2.3.
Objetivo de aprendizagem e desenvolvimento : ( EI02TS03 ) - Utilizar diferentes fontes sonoras disponíveis no ambiente em \textbf{brincadeiras} cantadas, canções, músicas e melodias.
Aprendizagens específicas - Ampliar repertório musical a partir de músicas e \textbf{brincadeiras} cantadas. - Desenvolver atitude de apreciação musical, aprimorando o gosto e a sensibilidade em relação à música. - Apreciar produções audiovisuais como musicais, \textbf{brinquedos} cantados, teatros de fantoches e marionetes. - Desenvolver o gosto por variados estilos musicais, reproduzindo -os nas \textbf{brincadeiras} cantadas e em suas criações musicais. - Desenvolver a capacidade de \textbf{escuta} dos \textbf{sons}, identificando -os e correlacionando -os aos objetos ou seres que os produzem. - Ampliar a capacidade de reproduzir os \textbf{sons}, \textbf{imitando} -os de diversas formas ( \textbf{batendo} \textbf{palmas}, dançando, cantando, utilizando objetos sonoros, entre outras ). - Ampliar seu conhecimento de mundo por meio da investigação e descoberta das diferentes fontes sonoras ( natureza, instrumentos musicais, equipamentos audiovisuais, corpo, etc. ). - Reconhecer \textbf{sons} familiares e identificar possibilidades de reproduzi -los utilizando objetos de seu cotidiano ou instrumentos musicais. - Desenvolver a socialização e o sentimento de pertença, participando de jogos e \textbf{brincadeiras} em grupo, que envolvam a dança e a música.
Sugestões de experiências - \textbf{Imitar} \textbf{sons} da natureza – canto dos pássaros, “ vozes ” de animais, \textbf{sons} do vento, da chuva e outros. - \textbf{Imitar} \textbf{sons} da cultura – vozes humanas, \textbf{sons} de instrumentos musicais e de máquinas. - Participar de eventos sonoros e produções musicais. - Participar de jogos e \textbf{brincadeiras} que envolvam a dança e/ou improvisação musical. - Explorar objetos de seu cotidiano ou instrumentos musicais. - Reproduzir \textbf{sons} ou canções conhecidas em \textbf{brincadeiras}. - Apreciar canções e músicas de diferentes culturas. - Escutar músicas de diferentes tradições. - Cantar e \textbf{brincar} com música de diferentes gêneros e gostos. - \textbf{Imitar} e \textbf{inventar} música de sua preferência. - Ouvir os \textbf{sons} do próprio corpo e dos \textbf{colegas}. - Ouvir os mais variados \textbf{sons} do entorno e estar atento ao silêncio. - Ouvir músicas de diferentes estilos. - Ouvir histórias em que se possa reproduzir os \textbf{sons} presentes no enredo. - Cantar com \textbf{acompanhamento} de \textbf{sons} produzidos por batidas ritmadas em várias partes do corpo ( mãos, \textbf{pés}, \textbf{dedos}, barriga, etc. ). - \textbf{Brincar} com jogos cantados e rítmicos. - Utilizar diferentes objetos e deles extrair \textbf{som} para acompanhar canções conhecidas. - Explorar por meio de \textbf{brincadeiras} diversas fontes sonoras, discriminando -as. 19.3.
Avaliação - \textbf{acompanhamento} do processo de aprendizagem e desenvolvimento das \textbf{crianças} bem \textbf{pequenas} : Observação \textbf{criteriosa}, \textbf{registro} e análise quanto : - À capacidade de criar \textbf{sons} com materiais, objetos e instrumentos musicais. - Ao desenvolvimento da sensibilidade, da capacidade de \textbf{manipular} materiais diversos e a criatividade. - À iniciativa em buscar e utilizar diferentes fontes sonoras para produzir \textbf{sons}.

% matched lemmas: acompanhamento, adulto, areia, autoconfiança, bater, brincadeira, brincar, brinquedo, colega, criança, criativamente, criterioso, cuidado, curiosidade, dedo, demonstrar, desejo, emitir, empilhar, escuta, gesto, imitar, infantil, inventar, lúdico, manipular, manusear, manuseio, montar, mês, mímico, palma, pequeno, peso, pintura, plástico, progressivamente, pé, registro, reutilizável, som, textura, tinta
\end{textsample}
